\section{Conclusion}
\label{sec:conclusion}

In this Phase~2 investigation we compared deep and hybrid models against the Phase~1 tree baselines on Tox21 under one frozen-configuration protocol (Phase~1 scaffold and random splits, three seeds). The evidence supports four conclusions, with the limits stated:
\begin{enumerate}
    \item \textbf{Tree models remain best on scaffold.} Random Forest ($0.4035$ PR~AUC) and XGBoost ($0.4010$) on Morgan fingerprints plus descriptors score above every deep model (H1 $0.3553$, FP-MLP $0.3429$, GCN $0.3397$, MPNN $0.3069$) and above both H2 variants.
    \item \textbf{Among the deep models.} The GCN scores higher than the plain MPNN ($11$ of $12$ endpoints, $p = 0.0015$). H1 has a higher observed mean PR~AUC than the FP-MLP ($+0.0124$, $9$ of $12$ endpoints), but the difference is not significant ($p = 0.151$).
    \item \textbf{Descriptor alignment matters for H1.} Permuting the descriptor rows across molecules lowers H1 from $0.3553$ to $0.2547$ PR~AUC, so its benefit depends on correct molecule--descriptor alignment. The experiment does not identify the mechanism.
    \item \textbf{Equal-weight rank fusion did not beat XGBoost on scaffold.} H2 (XGBoost plus the FP-MLP, selected on validation PR~AUC; $0.3945$) beats both of its deep components but not XGBoost alone ($4$ wins and $8$ losses, $p = 0.5186$). On the random split it is nominally ahead ($0.5727$ vs.\ $0.5647$) without significance ($p = 0.0923$).
\end{enumerate}

These conclusions rest on one dataset, one fixed split of each type, three seeds and 12 endpoints, so the tests have low power. In this benchmark, descriptor-based tree models remain the strongest single models, and the deep models and the ensemble did not surpass them on the scaffold split.
